\documentclass[12pt,a4paper]{article}

% ── Codificación y lenguaje ──────────────────────────────────────────────────
\usepackage[utf8]{inputenc}
\usepackage[T1]{fontenc}
\usepackage[spanish,mexico]{babel}

% ── Fuentes ──────────────────────────────────────────────────────────────────
\usepackage{lmodern}

% ── Márgenes ─────────────────────────────────────────────────────────────────
\usepackage[top=2.5cm, bottom=2.5cm, left=3cm, right=2.5cm]{geometry}

% ── Color ────────────────────────────────────────────────────────────────────
\usepackage[table]{xcolor}
\definecolor{rosado}{RGB}{220, 180, 185}
\definecolor{grisclaro}{RGB}{245, 245, 245}
\definecolor{grisoscuro}{RGB}{60, 60, 60}
\definecolor{azul}{RGB}{30, 70, 120}
\definecolor{verde}{RGB}{39, 130, 80}

% ── Gráficos ─────────────────────────────────────────────────────────────────
\usepackage{graphicx}
\usepackage{tikz}
\usepackage{pgfplots}
\pgfplotsset{compat=1.18}
\usetikzlibrary{calc, positioning, shapes.geometric}

% ── Tablas ───────────────────────────────────────────────────────────────────
\usepackage{booktabs}
\usepackage{tabularx}
\usepackage{longtable}
\usepackage{multirow}
\usepackage{array}
\usepackage{colortbl}

% ── Código ───────────────────────────────────────────────────────────────────
\usepackage{listings}
\lstset{
  basicstyle=\ttfamily\small,
  backgroundcolor=\color{grisclaro},
  frame=single,
  rulecolor=\color{azul!50},
  breaklines=true,
  showstringspaces=false,
  numbers=left,
  numberstyle=\tiny\color{grisoscuro},
  xleftmargin=12pt,
  aboveskip=8pt,
  belowskip=8pt,
  language=SQL,
  keywordstyle=\color{azul}\bfseries,
  stringstyle=\color{verde},
  commentstyle=\color{grisoscuro}\itshape,
}

\lstdefinelanguage{json}{
  basicstyle=\ttfamily\small,
  morestring=[b]",
  morekeywords={true, false, null},
  keywordstyle=\color{azul}\bfseries,
  stringstyle=\color{verde},
}

% ── Encabezados ──────────────────────────────────────────────────────────────
\usepackage{fancyhdr}
\pagestyle{fancy}
\fancyhf{}
\fancyhead[L]{\small\color{grisoscuro}AOS --- Primer Avance}
\fancyhead[R]{\small\color{grisoscuro}Sistema de Reservaciones para Restaurante}
\fancyfoot[C]{\small\thepage}
\renewcommand{\headrulewidth}{0.4pt}

% ── Hipervínculos ────────────────────────────────────────────────────────────
\usepackage[hidelinks]{hyperref}

% ── Secciones ────────────────────────────────────────────────────────────────
\usepackage{titlesec}
\titleformat{\section}{\large\bfseries}{}{0em}{\thesection.\quad}[\vspace{2pt}\hrule\vspace{4pt}]
\titleformat{\subsection}{\normalsize\bfseries}{\thesubsection.}{0.5em}{}
\titleformat{\subsubsection}{\normalsize\bfseries\itshape}{\thesubsubsection.}{0.5em}{}

% ── Otros ────────────────────────────────────────────────────────────────────
\usepackage{enumitem}
\usepackage{float}
\usepackage{caption}
\captionsetup{font=small, labelfont=bf}

% ════════════════════════════════════════════════════════════════════════════
\begin{document}
\pagenumbering{gobble}  % Sin número en portada

% ══════════════════════════════════════════════════════
%  PORTADA — ESTILO UNIVERSIDAD POLITÉCNICA DE PACHUCA
% ══════════════════════════════════════════════════════
\begin{titlepage}
\newgeometry{top=2cm, bottom=2cm, left=3.5cm, right=3.5cm}

% Triángulo decorativo rosa en esquina superior derecha
\begin{tikzpicture}[remember picture, overlay]
  \fill[rosado!70]
    (current page.north east)
    -- ++(-3.5cm, 0)
    -- ++(0, -4.5cm)
    -- cycle;
\end{tikzpicture}

\vspace*{1.5cm}

% Universidad
\begin{center}
  {\small\scshape Universidad Politécnica de Pachuca}
\end{center}

\vspace{2.5cm}

% Título principal
\begin{center}
  {\LARGE\bfseries ARQUITECTURA\\[4pt]
  ORIENTADA A\\[4pt]
  SERVICIOS}
\end{center}

\vspace{1cm}

% Subtítulo
\begin{center}
  {\large\bfseries PRIMER AVANCE: SISTEMA DE\\[4pt]
  RESERVACIONES PARA RESTAURANTE}
\end{center}

\vspace{2.2cm}

% Datos del proyecto
\noindent\textbf{DOCENTE:} Jazmín Rodrigues Flores\\[6pt]
\noindent\textbf{ALUMNOS:}\\[4pt]
\hspace*{1em} Jovanny Hernández Hernández\\[2pt]
\hspace*{1em} Cristopher López Suárez\\[6pt]
\noindent\textbf{ASIGNATURA:} Arquitectura Orientada a Servicios\\[6pt]
\noindent\textbf{GRUPO:} SFTW\_09\_01\\[6pt]
\noindent\textbf{FECHA:} Julio de 2026

\vfill

\begin{center}
  {\small Ingeniería en Software}
\end{center}

\restoregeometry
\end{titlepage}

% ── Activar numeración y encabezados desde aquí ──────────────────────────────
\pagenumbering{arabic}
\setcounter{page}{1}

% ── ÍNDICE ───────────────────────────────────────────────────────────────────
\tableofcontents
\newpage

% ════════════════════════════════════════════════════════════════════════════
\section{Introducción}

El presente documento corresponde al \textbf{Primer Avance} del proyecto
\textit{Sistema de Reservaciones para Restaurante}, desarrollado en el marco
de la asignatura \textbf{Arquitectura Orientada a Servicios (AOS)} de la
carrera de Ingeniería en Software de la Universidad Politécnica de Pachuca.

El objetivo del sistema es digitalizar el proceso de reserva de mesas,
gestión del catálogo de platillos y atención al cliente en un entorno
restaurantero, utilizando una arquitectura que combina base de datos
relacional (MySQL) y base de datos NoSQL (MongoDB).



% ════════════════════════════════════════════════════════════════════════════
\section{Casos de Uso (CU), Historias de Usuario (HU) y Requisitos}

\subsection{Actores del Sistema}

\begin{center}
\begin{tabular}{lp{9cm}}
\toprule
\rowcolor{grisclaro}
\textbf{Actor} & \textbf{Descripción} \\
\midrule
Cliente       & Usuario que consulta el menú y realiza reservaciones. \\
Mesero        & Gestiona mesas y consulta las reservaciones del día. \\
Administrador & Administra el sistema completo: usuarios, menú y reportes. \\
\bottomrule
\end{tabular}
\end{center}

\subsection{Casos de Uso Principales}

\begin{center}
{\small
\begin{tabular}{clp{6cm}l}
\toprule
\rowcolor{grisclaro}
\textbf{CU\#} & \textbf{Nombre} & \textbf{Descripción} & \textbf{Actor} \\
\midrule
CU-01 & Registrar usuario     & El cliente crea una cuenta con nombre, email y contraseña. & Cliente \\
\rowcolor{grisclaro}
CU-02 & Iniciar sesión        & Acceso al sistema con credenciales válidas. & Todos \\
CU-03 & Ver menú              & Consulta de platillos con precios. & Cliente \\
\rowcolor{grisclaro}
CU-04 & Realizar reservación  & Selección de fecha, hora, mesa y personas. & Cliente \\
CU-05 & Cancelar reservación  & Cancelación de una reservación activa. & Cliente \\
\rowcolor{grisclaro}
CU-06 & Ver mis reservaciones & Historial de reservaciones del usuario. & Cliente \\
CU-07 & Gestionar mesas       & Alta, baja y modificación de mesas. & Admin/Mesero \\
\rowcolor{grisclaro}
CU-08 & Gestionar menú        & CRUD del catálogo de platillos. & Admin \\
CU-09 & Ver reservaciones del día & Listado de reservaciones en curso. & Admin/Mesero \\
\rowcolor{grisclaro}
CU-10 & Generar reporte       & Reporte de ocupación por fechas. & Admin \\
\bottomrule
\end{tabular}}
\end{center}

\subsection{Historias de Usuario con Puntos de Historia (Story Points)}

\begin{center}
{\small
\begin{tabular}{cp{7.5cm}ccc}
\toprule
\rowcolor{grisclaro}
\textbf{HU\#} & \textbf{Historia de Usuario} & \textbf{Prior.} & \textbf{Pts.} & \textbf{Spr.} \\
\midrule
HU-01 & \textbf{Como} cliente, \textbf{quiero} registrarme \textbf{para} hacer reservaciones fácilmente. & Alta & 3 & 1 \\
\rowcolor{grisclaro}
HU-02 & \textbf{Como} cliente, \textbf{quiero} ver el menú con precios \textbf{para} elegir antes de llegar. & Alta & 5 & 1 \\
HU-03 & \textbf{Como} cliente, \textbf{quiero} reservar fecha, hora y personas \textbf{para} asegurar mi lugar. & Alta & 8 & 1 \\
\rowcolor{grisclaro}
HU-04 & \textbf{Como} cliente, \textbf{quiero} recibir confirmación de mi reservación. & Alta & 3 & 1 \\
HU-05 & \textbf{Como} cliente, \textbf{quiero} cancelar mi reservación \textbf{para} liberar el lugar. & Media & 5 & 2 \\
\rowcolor{grisclaro}
HU-06 & \textbf{Como} cliente, \textbf{quiero} ver mi historial de reservaciones. & Media & 3 & 2 \\
HU-07 & \textbf{Como} admin, \textbf{quiero} gestionar el menú \textbf{para} mantenerlo actualizado. & Alta & 8 & 2 \\
\rowcolor{grisclaro}
HU-08 & \textbf{Como} admin, \textbf{quiero} ver reservaciones del día \textbf{para} organizar al personal. & Media & 5 & 2 \\
HU-09 & \textbf{Como} mesero, \textbf{quiero} ver el estado de mesas \textbf{para} organizarlas. & Media & 5 & 3 \\
\rowcolor{grisclaro}
HU-10 & \textbf{Como} admin, \textbf{quiero} generar reportes de ocupación \textbf{para} tomar decisiones. & Baja & 8 & 3 \\
\midrule
\multicolumn{3}{r}{\textbf{Total}} & \textbf{53} & --- \\
\bottomrule
\end{tabular}}
\end{center}

\subsection{Requisitos Funcionales}

\begin{enumerate}[label=\textbf{RF-\arabic*:}, leftmargin=2.8cm]
  \item El sistema permitirá registrar, editar y eliminar usuarios.
  \item El sistema permitirá iniciar y cerrar sesión con validación de credenciales.
  \item El sistema mostrará el catálogo de platillos con nombre, descripción, precio y categoría.
  \item El sistema permitirá crear reservaciones con fecha, hora, mesa y número de personas.
  \item El sistema validará que no existan reservaciones duplicadas en la misma mesa y horario.
  \item El sistema permitirá cancelar reservaciones con al menos 1 hora de anticipación.
  \item El sistema mostrará el historial de reservaciones por cliente.
  \item El sistema permitirá al administrador gestionar el menú de productos (CRUD).
  \item El sistema mostrará un tablero de mesas con disponibilidad en tiempo real.
  \item El sistema generará reportes de reservaciones filtrados por fecha.
\end{enumerate}

\subsection{Requisitos No Funcionales}

\begin{enumerate}[label=\textbf{RNF-\arabic*:}, leftmargin=2.8cm]
  \item El sistema deberá responder en menos de 2 segundos por petición.
  \item Las contraseñas deberán almacenarse con hash (bcrypt).
  \item El sistema deberá ser \textit{responsive} (adaptable a móvil y escritorio).
  \item El sistema deberá estar disponible el 99\% del tiempo durante horario de operación.
  \item Solo usuarios autenticados podrán realizar reservaciones (JWT).
  \item El sistema deberá soportar al menos 50 usuarios concurrentes.
\end{enumerate}

% ════════════════════════════════════════════════════════════════════════════
\section{Base de Datos NoSQL --- Catálogo de Productos (MongoDB)}

Se utiliza \textbf{MongoDB} para el catálogo de platillos debido a la naturaleza flexible
y variable de los datos: un platillo puede tener opciones extras, alérgenos e ingredientes
variables. Esta estructura dinámica se adapta mejor a documentos JSON que a tablas relacionales rígidas.

\subsection{Colección: \texttt{menu\_items}}

\begin{lstlisting}[language=json, caption={Documento JSON --- Colección menu\_items}]
{
  "_id": "ObjectId('665abc123def')",
  "nombre": "Tacos de Res",
  "descripcion": "3 tacos de res asada con cebolla, cilantro y salsa",
  "categoria": "Platillo principal",
  "precio": 89.00,
  "imagen_url": "https://restaurante.com/img/tacos-res.jpg",
  "disponible": true,
  "ingredientes": ["res", "cebolla", "cilantro", "tortilla"],
  "alergenos": ["gluten"],
  "opciones": [
    { "nombre": "Sin cebolla",   "costo_extra": 0 },
    { "nombre": "Extra queso",   "costo_extra": 15.00 }
  ],
  "tiempo_preparacion_min": 15,
  "calorias": 450,
  "tags": ["popular", "picante"],
  "fecha_creacion": "2025-01-15T10:00:00Z"
}
\end{lstlisting}

\subsection{Colección: \texttt{mesas}}

\begin{lstlisting}[language=json, caption={Documento JSON --- Colección mesas}]
{
  "_id": "ObjectId('665def456abc')",
  "numero": 5,
  "capacidad": 4,
  "ubicacion": "Terraza",
  "disponible": true,
  "caracteristicas": ["vista al jardin", "techada"],
  "activa": true
}
\end{lstlisting}

% ════════════════════════════════════════════════════════════════════════════
\section{Base de Datos Relacional --- Esquema Completo (MySQL)}

Se utiliza \textbf{MySQL} para los datos estructurados y críticos del negocio: usuarios,
reservaciones, estados y reseñas. Estos datos requieren integridad referencial,
transacciones ACID y relaciones bien definidas entre entidades.

\subsection{Diagrama Entidad-Relación}

\begin{figure}[H]
\centering
\begin{tikzpicture}[
  ent/.style={draw, rectangle, rounded corners=3pt, minimum width=2.8cm,
              minimum height=0.8cm, font=\bfseries\small, fill=grisclaro},
  rel/.style={draw, diamond, aspect=2, minimum width=2cm,
              minimum height=0.8cm, font=\small\itshape, fill=white},
  lbl/.style={font=\footnotesize},
  node distance=0.8cm and 2.6cm
]
  \node[ent] (U)  {USUARIOS};
  \node[ent] (R)  [right=5cm of U] {RESERVACIONES};
  \node[ent] (M)  [right=5cm of R] {MESAS};
  \node[ent] (E)  [below=2cm of R] {ESTADOS};
  \node[ent] (Re) [below=2cm of U] {RESE\~NAS};

  % Relacion USUARIOS -- RESERVACIONES
  \node[rel] (hace)    at ($(U)!0.5!(R)$)   {hace};
  \draw[-] (U) -- node[above,lbl]{1} (hace);
  \draw[-] (hace) -- node[above,lbl]{N} (R);

  % Relacion RESERVACIONES -- MESAS
  \node[rel] (asigna)  at ($(R)!0.5!(M)$)   {asigna};
  \draw[-] (R) -- node[above,lbl]{N} (asigna);
  \draw[-] (asigna) -- node[above,lbl]{1} (M);

  % Relacion RESERVACIONES -- ESTADOS
  \draw[-] (R) -- node[right,lbl]{N} (E);

  % Relaciones con RESENAS
  \draw[dashed,-] (U) -- node[left,lbl]{1} (Re);
  \draw[dashed,-] (R) -- node[below,lbl]{1} (Re);

\end{tikzpicture}
\caption{Diagrama Entidad-Relación --- Sistema de Reservaciones}
\end{figure}

\subsection{Tablas del Esquema}

\subsubsection{Tabla \texttt{usuarios}}
\begin{lstlisting}[caption={SQL --- Tabla usuarios}]
CREATE TABLE usuarios (
    id_usuario     INT PRIMARY KEY AUTO_INCREMENT,
    nombre         VARCHAR(100) NOT NULL,
    apellido       VARCHAR(100) NOT NULL,
    email          VARCHAR(150) UNIQUE NOT NULL,
    telefono       VARCHAR(15),
    password_hash  VARCHAR(255) NOT NULL,
    rol   ENUM('cliente','mesero','admin') DEFAULT 'cliente',
    activo         BOOLEAN DEFAULT TRUE,
    fecha_registro DATETIME DEFAULT CURRENT_TIMESTAMP
);
\end{lstlisting}

\subsubsection{Tabla \texttt{mesas}}
\begin{lstlisting}[caption={SQL --- Tabla mesas}]
CREATE TABLE mesas (
    id_mesa   INT PRIMARY KEY AUTO_INCREMENT,
    numero    INT UNIQUE NOT NULL,
    capacidad INT NOT NULL,
    ubicacion VARCHAR(100),
    activa    BOOLEAN DEFAULT TRUE
);
\end{lstlisting}

\subsubsection{Tabla \texttt{estados\_reservacion}}
\begin{lstlisting}[caption={SQL --- Tabla estados\_reservacion}]
CREATE TABLE estados_reservacion (
    id_estado INT PRIMARY KEY AUTO_INCREMENT,
    nombre    VARCHAR(50) NOT NULL
    -- Valores: Pendiente, Confirmada, Cancelada, Completada
);
\end{lstlisting}

\subsubsection{Tabla \texttt{reservaciones}}
\begin{lstlisting}[caption={SQL --- Tabla reservaciones (tabla central)}]
CREATE TABLE reservaciones (
    id_reservacion INT PRIMARY KEY AUTO_INCREMENT,
    id_usuario     INT NOT NULL,
    id_mesa        INT NOT NULL,
    id_estado      INT NOT NULL DEFAULT 1,
    fecha          DATE NOT NULL,
    hora_inicio    TIME NOT NULL,
    hora_fin       TIME NOT NULL,
    num_personas   INT NOT NULL,
    notas          TEXT,
    fecha_creacion DATETIME DEFAULT CURRENT_TIMESTAMP,
    FOREIGN KEY (id_usuario) REFERENCES usuarios(id_usuario),
    FOREIGN KEY (id_mesa)    REFERENCES mesas(id_mesa),
    FOREIGN KEY (id_estado)  REFERENCES estados_reservacion(id_estado),
    CONSTRAINT chk_horario CHECK (hora_fin > hora_inicio),
    UNIQUE KEY uk_mesa_fecha_hora (id_mesa, fecha, hora_inicio)
);
\end{lstlisting}

\subsubsection{Tabla \texttt{resenas}}
\begin{lstlisting}[caption={SQL --- Tabla resenas}]
CREATE TABLE resenas (
    id_resena      INT PRIMARY KEY AUTO_INCREMENT,
    id_usuario     INT NOT NULL,
    id_reservacion INT NOT NULL,
    calificacion   TINYINT NOT NULL CHECK (calificacion BETWEEN 1 AND 5),
    comentario     TEXT,
    fecha          DATETIME DEFAULT CURRENT_TIMESTAMP,
    FOREIGN KEY (id_usuario)     REFERENCES usuarios(id_usuario),
    FOREIGN KEY (id_reservacion) REFERENCES reservaciones(id_reservacion)
);
\end{lstlisting}

\subsubsection{Tabla \texttt{menu\_relacional} (referencia a MongoDB)}
\begin{lstlisting}[caption={SQL --- Tabla menu\_relacional}]
CREATE TABLE menu_relacional (
    id_item    INT PRIMARY KEY AUTO_INCREMENT,
    mongo_id   VARCHAR(50) UNIQUE,   -- Referencia al _id de MongoDB
    nombre     VARCHAR(150) NOT NULL,
    categoria  VARCHAR(100),
    precio     DECIMAL(10,2) NOT NULL,
    disponible BOOLEAN DEFAULT TRUE
);
\end{lstlisting}

% ════════════════════════════════════════════════════════════════════════════
\section{Metodología --- Scrum}



\subsection{Product Backlog con HUs Valoradas}

\begin{center}
{\small
\begin{tabular}{cp{7cm}ccc}
\toprule
\rowcolor{grisclaro}
\textbf{HU\#} & \textbf{Historia de Usuario} & \textbf{Prior.} & \textbf{Pts.} & \textbf{Spr.} \\
\midrule
HU-01 & Registrar usuario              & Alta  & 3  & 1 \\
\rowcolor{grisclaro}
HU-02 & Ver menú con precios           & Alta  & 5  & 1 \\
HU-03 & Realizar reservación           & Alta  & 8  & 1 \\
\rowcolor{grisclaro}
HU-04 & Recibir confirmación           & Alta  & 3  & 1 \\
HU-05 & Cancelar reservación           & Media & 5  & 2 \\
\rowcolor{grisclaro}
HU-06 & Ver historial de reservaciones & Media & 3  & 2 \\
HU-07 & Gestionar menú (Admin)         & Alta  & 8  & 2 \\
\rowcolor{grisclaro}
HU-08 & Ver reservaciones del día      & Media & 5  & 2 \\
HU-09 & Tablero de estado de mesas     & Media & 5  & 3 \\
\rowcolor{grisclaro}
HU-10 & Generar reportes de ocupación  & Baja  & 8  & 3 \\
\midrule
\multicolumn{3}{r}{\textbf{Total}} & \textbf{53} & --- \\
\bottomrule
\end{tabular}}
\end{center}

\subsection{Planificación de Sprints}

\begin{center}
\begin{tabular}{cccc}
\toprule
\rowcolor{grisclaro}
\textbf{Sprint} & \textbf{Duración} & \textbf{HUs} & \textbf{Puntos} \\
\midrule
Sprint 1 & 1 semana & HU-01, HU-02, HU-04            & 11 pts \\
Sprint 2 & 1 semana & HU-03                           &  8 pts \\
Sprint 3 & 1 semana & HU-05, HU-06, HU-08             & 13 pts \\
Sprint 4 & 1 semana & HU-07, HU-09                    & 13 pts \\
Sprint 5 & 1 semana & HU-10                           &  8 pts \\
\midrule
\multicolumn{3}{r}{\textbf{Total}} & \textbf{53 pts} \\
\bottomrule
\end{tabular}
\end{center}

\subsection{Gráfica de Quemado (Burndown Chart)}

\begin{figure}[H]
\centering
\begin{tikzpicture}
\begin{axis}[
  width=13cm, height=7cm,
  xlabel={\textbf{Semanas}},
  ylabel={\textbf{Puntos Restantes}},
  xmin=0, xmax=5,
  ymin=0, ymax=58,
  xtick={0,1,2,3,4,5},
  xticklabels={Inicio, Sem.1, Sem.2, Sem.3, Sem.4, Sem.5},
  ytick={0,10,20,30,40,50,53},
  grid=both,
  grid style={dotted, gray!40},
  legend pos=north east,
  title={\textbf{Burndown Chart --- Sistema de Reservaciones}},
]
% Línea ideal
\addplot[blue!70!black, thick, dashed, mark=square*,
         mark options={fill=blue!70!black}]
  coordinates {(0,53)(1,42)(2,32)(3,21)(4,11)(5,0)};
\addlegendentry{Quemado ideal}

% Punto de inicio real
\addplot[red!70!black, thick, mark=*, mark options={fill=red!70!black}]
  coordinates {(0,53)};
\addlegendentry{Quemado real}

% Líneas divisoras de sprint (cada semana es un sprint)
\draw[dashed, gray] (axis cs:1,0) -- (axis cs:1,55);
\draw[dashed, gray] (axis cs:2,0) -- (axis cs:2,55);
\draw[dashed, gray] (axis cs:3,0) -- (axis cs:3,55);
\draw[dashed, gray] (axis cs:4,0) -- (axis cs:4,55);
\node[font=\footnotesize] at (axis cs:0.5,56) {S1};
\node[font=\footnotesize] at (axis cs:1.5,56) {S2};
\node[font=\footnotesize] at (axis cs:2.5,56) {S3};
\node[font=\footnotesize] at (axis cs:3.5,56) {S4};
\node[font=\footnotesize] at (axis cs:4.5,56) {S5};
\end{axis}
\end{tikzpicture}
\caption{Gráfica de Quemado --- Velocidad ideal por sprint (53 puntos / 5 semanas)}
\end{figure}

\begin{center}
\begin{tabular}{ccc}
\toprule
\rowcolor{grisclaro}
\textbf{Semana} & \textbf{Puntos Ideales Restantes} & \textbf{Puntos Reales Restantes} \\
\midrule
0 (Inicio) & 53 & 53 \\
1 (Fin S1) & 42 & --- \\
2 (Fin S2) & 32 & --- \\
3 (Fin S3) & 21 & --- \\
4 (Fin S4) & 11 & --- \\
5 (Fin S5) &  0 & --- \\
\bottomrule
\end{tabular}
\end{center}

\subsection{Herramientas del Proyecto}

\begin{center}
{\small
\begin{tabular}{llp{6cm}}
\toprule
\rowcolor{grisclaro}
\textbf{Herramienta} & \textbf{Uso} & \textbf{Descripción} \\
\midrule
GitHub        & Control de versiones & Repositorio del código fuente. \\
Trello        & Gestión ágil         & Tablero Scrum: Backlog / En progreso / Hecho. \\
Google Drive  & Documentación        & Archivos y documentos compartidos. \\
MongoDB Atlas & BD NoSQL             & Catálogo de platillos (menú). \\
MySQL         & BD Relacional        & Usuarios, mesas y reservaciones. \\
Docker        & Contenedores         & Entorno de desarrollo consistente. \\
\bottomrule
\end{tabular}}
\end{center}

% ════════════════════════════════════════════════════════════════════════════
\section{Conclusiones}

En este Primer Avance se establecieron las bases del \textbf{Sistema de Reservaciones
para Restaurante}, cubriendo los siguientes entregables:

\begin{itemize}[label={$\bullet$}]
  \item Definición de actores, 10 Casos de Uso y 10 Historias de Usuario valoradas bajo el marco Scrum.
  \item Levantamiento de 10 requisitos funcionales y 6 no funcionales del sistema.
  \item Diseño de BD NoSQL (MongoDB) para el catálogo de productos con estructura flexible.
  \item Diseño completo de BD Relacional (MySQL) con integridad referencial y constraints.
  \item Planificación de la metodología Scrum: Product Backlog, 5 sprints y Burndown Chart.
  \item Identificación del stack tecnológico y herramientas de trabajo colaborativo.
\end{itemize}

El proyecto tiene una duración estimada de \textbf{5 semanas} en 5 sprints de 1 semana cada uno,
con un total de \textbf{53 puntos de historia}, desarrollado por
\textbf{Jovanny Hernández Hernández} y \textbf{Cristopher López Suárez}.

\end{document}
